\documentclass[../../main.tex]{subfiles}

\begin{document}
\chapter{Temporal-Difference Learning}
\makeSubfileHeader{Temporal-Difference Learning}

\section{Introduction}
    \textbf{Temporal-Difference learning} (TD) is one of the central ideas in reinforcement learning.
    Like DP, TD methods update estimates based in part on other learned estimates. On the other hand, TD methods learn
    directly from experiences like MC methods. DP, MC, and TD methods' relationships will be kept founded in later chapters.
    n-step algorithms and TD($\lambda$) algorithms are examples of them.

    The main advantages of TD over MC and DP methods is as following:
    \begin{itemize}
        \item Do not have to wait until the epsidoes ends.
        \item Do not have to ignore or discount episodes on which experimental actions are taken. 
        \item No model is required.
    \end{itemize}

\section{TD Prediction}
    TD method's advantage over MC method is that it can be applied to continuing tasks, since it updates each time step.
    The simplest form of TD method makes the update
    \begin{equation}
        V(S_t) \leftarrow V(S_t) + \alpha \Big[ R_{t+1} + \gamma V(S_{t+1} - V(S_t))\Big] \label{eq:TDUpdate}
    \end{equation}
    immediately on transition to $S_t$ receiving $R_{t+1}$. This method is called \textbf{TD(0)}, and is a speacial case of n-step bootstrap and TD($\lambda$).

    \subsection{Algorithm}
        \Algorithm{tabularTDZeroPred}
        {Tabular TD(0) Prediction}
        {\inputitem{$\pi$}{policy to be evaluated} \\
        \inputitem{$\alpha \in (0,1]$}{step size parameter}}
        {\inputitem{$v_\pi$}{state-value function of policy $\pi$}}
        {
            \tcp{initialize}
            \ForEach{$s \in \mathcal{S}$}
            {
                $V(s) \leftarrow 
                \begin{cases}
                    \text{arbitrary value} &\text{if } s \in \mathcal{S^+} \\
                    0 &\text{if } s \in \mathcal{S} \setminus \mathcal{S^+}
                \end{cases}$\;
            }
            \ForEach{episodes}
            {
                $S \leftarrow \text{arbitrary } s \in \mathcal{S^+}$\;
                \ForEach{step in episode}
                {
                    $A \leftarrow$ action given by $\pi$ for $S$\;
                    take action $A$, receive reward $R$ and next state $S'$\;
                    $V(S) \leftarrow V(S) + \alpha\big[R + \gamma V(S') - V(S)\big]$\;
                    $S \leftarrow S'$\;
                    \If{$S \in \mathcal{S} \setminus \mathcal{S}^+$}
                    {
                        break\;
                    }
                }
            }
            \Return{$V \simeq v_\pi$}
        }

        Since TD(0) method updates value with existing estimates, we say that it \textbf{bootstraps}, like DP method.
        TD combines the sampling of MC and bootstrapping of DP, and we call this update as \textbf{sampling update}.
        Note that it is different from \emph{expected update} of DP.

        \subsubsection{Convergence of the Algorithm}
            Since TD(0) method uses estimates instead of complete $G_t$, one may be skeptical of its convergence.
            However, it is known to converge to $v_\pi$, happily. For any fixed policy $\pi$, TD(0) is known to converge to $v_\pi$
            with constant step-size parameter if it is sufficiently small, and is known to converge to $v_\pi$ with probability 1 if
            the step-size parameter changes step by step while satisfying the stochastic approximation conditions \ref{thm:condForConv}.

            There is no known answer to the question "If TD and MC both converges to correct answer, then which one converges more quickly?".
            In practice, TD methods usually converge faster than MC methods on stocahstic tasks.

    \subsection{TD Error}
        The following definition of \emph{TD error} arises in various forms throughout reinforcement learning:
        \begin{defn}{TD Error}{TDError}
            Let $V$ be the estimate of state-value function. Then, the \textbf{TD error} $\delta_t$ at time $t$ is defined as
            \begin{equation}
                \delta_t = R_{t+1} + \gamma V(S_{t+1}) - V(S_t)
            \end{equation}
        \end{defn}
        Note that TD error is also contained in algorithm \ref{alg:tabularTDZeroPred}. If the estimate $V$ does not change during the episode,
        the Monte Carlo error $G_t - V(S_t)$ can be expressed as the weighted sum of TD errors;
        \begin{equation}
            G_t - V(S_t) = \sum_{k=t}^{T-1} \gamma^{k-t}\delta_t.
        \end{equation}
        Even though $V$ changes during episode, the equation above may still hold approximately if the step size is small enough.
        Generalization of this idea plays an important role in the theory and algorithms of TD learning. 

    \subsection{Batch Updates on TD}
        We can apply batch updates to TD(0) methods by computing the increments \ref{eq:TDUpdate} every time step but updating the value
        after the batch, by the sum of all the increments.

        Very interesting property can be found on batch TD(0) prediction. Batch TD(0) prediction converges to the \textbf{certainty-equivalence estimate}
        of the MDP. Certainty-equivalence estimate of MDP assumes that the estimate of the MDP is perfect. That is, it is \emph{certain} about its estimate.

        To calculate certainty-equivalence estimate directly, it requires memory complexity of $\mathcal{O}(n^2)$ due to the transition probability table,
        and time complexity of $\mathcal{O}(n^3)$ due to solving the Bellman equation as the linear system of equations.
        However, using TD(0) method, it requires memory complexity of only $\mathcal{O}(n)$, to save each values of states, which is a surprising result.

        Meanwhile, batch MC method converges to minimizer of the mean-square error.
        
        \begin{rem}{TD(0) Converges to Certainty-equivalence Estimate}
            Suppose that the algorithm converged to fixed point. Then, the next batch update of value of each state $s$ should be 0.
            \begin{equation}
                \sum_{k \in \text{batch}}\sum_{t \in \{t : S_t^k=s\}}\alpha(R_{t+1}^k + \gamma V(S_{t+1}^k) - V(s))=0 \label{eq:rem1}
            \end{equation}
            Let $N(s)$ and $N(s,s')$ denote the number of visits of state $s$ in batch and the number of transition from $s$ to $s'$ in batch, respectively.
            Then, the equation \ref{eq:rem1} can be written as
            \begin{equation}
                \left(\sum_{t \in \{t : S_t = s\}} R_{t+1}\right) + \gamma \left( \sum_{s'}N(s,s')V(s') \right) - N(s)V(s) = 0. \label{eq:rem2}
            \end{equation}
            Dividing both side of equation \ref{eq:rem2} with $N(s)$, we get
            \begin{equation}
                \frac{1}{N(s)}\left( \sum_{t \in \{t : S_t = s\}} R_{t+1}\right) + \gamma \sum_{s'}\frac{N(s,s')}{N(s)}V(s') - V(s) = 0
            \end{equation}
            Then, the term $\hat{R}(s) = \frac{1}{N(s)}\left( \sum_{t \in \{t : S_t = s\}} R_{t+1}\right)$ is the average reward of state $s$ and
            the term $\hat{P}(s'|s) = \frac{N(s,s')}{N(s)}$ is the MLE estimate of transition probability from $s$ to $s'$.
            Substituting, we get
            \begin{equation}
                V(s) = \hat{R}(s) + \gamma \sum_{s'}\hat{P}(s'|s)V(s'),
            \end{equation}
            which equals the Bellman equation with certainty-equivalence estimate.
        \end{rem}

\section{TD Control}
    \subsection{Sarsa: On-policy TD Control}
    \subsection{Q-learning: Off-policy TD Control}

\end{document}